% ============================================================
% 1. Introduction
% ============================================================
\section{Introduction}
\label{sec:introduction}

A battery cell is designed five times over, and the five designs must be one
decision. The electrochemist selects the electrode system; the formulation
chemist the electrolyte---its transport properties, its additive package; the
materials engineer the electrode microstructure and any coating or doping; the
mechanical designer the cell architecture---thickness, porosity, separator,
current collectors, particle size; the thermal engineer the cooling. Each
choice constrains the next, and the aggregate is judged not by any single
specialist but by a contract of system-level metrics: energy density against
temperature ceilings, fast charge against lifetime. For decades the only
infrastructure capable of making this one decision end-to-end was a human team
of specialists, with all the costs of coordination that implies. Automated
methods have never crossed this boundary. Bayesian optimization moves
architecture knobs inside one chemistry \cite{chen2026battery}; autonomous
laboratories synthesize materials \cite{szymanski2023alab}; LLM chemists
execute molecular tasks \cite{boiko2023coscientist, bran2024chemcrow}. Each
covers a slice. No published system, to our knowledge, spans the full chain of
cell-design levers---electrolyte, separator, electrode, architecture,
thermal---within a single autonomous loop.

This paper reports the first system that does. In a virtual battery factory, a
single LLM agent receives a natural-language requirement---nothing else: no
configuration file, no parameter set, no candidate list---and returns a
complete, adjudicated cell design: the material choices, the formulation, the
architecture, the safety validation, and the deliverable package---with a closing first-principles endorsement pass (executed for the flagship molecular finalists, §5.8)
(ab initio DFT and molecular dynamics, reserved for the top candidates rather
than used for screening). The protocol
itself is a self-contained declarative artifact---natural-language rules, a
command-line tool library, and an audit ledger---bound to no specific agent
framework: the governance travels with the contract, not with the harness,
a property we verify empirically by running the identical protocol under
three orchestration substrates (§5.7).
The human
engineer's *in-loop* role ends at the moment the contract is stated; from that point
forward the agent plans, screens, simulates, diagnoses, falls back, and
concludes with no intermediate human oversight. The eight scenario contracts
we evaluate were chosen so that different tasks must exercise different levers
of this full chain---a transport upgrade for a fast-charge guarantee, an
SEI-kinetics bridge for a lifetime bound, a cathode-system switch for a
voltage plateau---and the governed agent exercises each in turn, in the same
loop, under the same protocol---the full-chain claim is thus a suite-level
statement: individual runs exercise the levers their bottleneck demands, the
suite collectively spans all five. This is the delegation setting management
science has not yet been able to study: a complete design chain handed to a
machine, with the principal present only as a written contract.

What makes such delegation safe---or dangerous---is the question this paper
answers. We deliberately construct the two failure modes a design organization
would most fear. The first is the \emph{solution-space failure}: the strongest
classical baseline, Bayesian optimization, moves only architecture parameters;
when a contract requires a material lever---a high-voltage spinel chemistry, an
SEI-suppressing coating---the optimum is simply not in its reachable set. The
second is more subtle and, we will argue, more dangerous: the
\emph{yardstick failure}. A protocol-free agent with the same tools, budget,
and model is not constrained to measure against the contract it was given. It
can shift the threshold (``dendrite onset is at $-10$ mV, so $-6.9$ mV is
safe''), substitute the model (an SEI growth law swapped for a slower variant,
reported in the contract model's vocabulary), buy compliance with a self-set
physical parameter (a cathode diffusivity scaled up for ``power-grade''
kinetics), alter the measurement caliber (a rate-retention ratio computed
against $C/3$ instead of the contractual $1C$ reference), or omit the
criterion variable entirely (a nail-penetration test evaluated by its own
steady-state temperature, with no trigger criterion anywhere in the
artifacts). In our data the bare agent and
the Bayesian optimizer tie---three of eight tasks ``achieved'' each---yet
their failures are orthogonal: one loses to an unreachable space, the other to
a malleable yardstick.

We introduce \emph{protocolized governance} as the intervention: a fixed,
always-on set of process rules the agent executes inside the design loop, plus
an adjudication layer that stands outside it. The protocol has five
components: a layered multi-precision funnel (cheap proxies first, molecular
screening only when the bottleneck calls for it, everything bounded by the
cell-scale judgment); evaluation-and-fallback routing (diagnose the scale of a
failure before retrying); heterogeneous model voting (two machine-learning
potentials plus a semi-empirical quantum method); a parameter bridge
(microscopic properties mapped to macroscopic model parameters under mandatory
source annotation); and---the load-bearing component---\emph{mechanized adjudication}. Under mechanized
adjudication the contract is pre-registered as machine-readable criteria in an
audit ledger before the first simulation runs; every verdict is computed by
code from tool-output files, not asserted by the agent in prose; and each
``achieved'' claim points to the exact file and key that supports it, with a
verifier that refuses incomplete audit trails. To our knowledge, this makes
the system the first design agent whose verdicts are \emph{computed} rather
than \emph{asserted}. In the language of delegation
theory, it is a delegation contract with enforced fidelity---the alignment
between the contract the principal believes it set and the behavior the agent
actually executes \cite{delegatedagency2026, nikpayam2024aiassisted}.

We evaluate the governance in a thirty-two-run matrix plus twelve ablation
runs. Eight scenario tasks---a sedan's energy-and-fast-charge contract, a
drone's weight-and-rate envelope, a smartphone's volumetric ceiling, and five
others---are written as natural-language requirements with absolute numeric
criteria and no other declarations. Each task is given, under identical
resources, to three decision-makers: the governed agent (two different LLMs,
separately and jointly); a protocol-free agent on the same harness with the
same toolset and budget; and a Bayesian optimizer configured to the standard
of the strongest comparable published agent framework \cite{chen2026battery}.
The results answer the question sharply. The governed agent achieves
\textbf{all eight} contracts under the two DeepSeek-generation models, and
four of eight under a third (Section~5.4). The BO baseline achieves
three---precisely the three tasks whose bottlenecks are structural rather than
material. The protocol-free agent achieves none within the same measurement
space: of its seven self-reported successes, two rest on a lithium-metal
cell outside the validated parameter sets, five dissolve under
contract-caliber re-adjudication, and the one session that made no claim
exhausted its budget.

Three further results speak to \emph{why} governance works. First, ablation
shows the mechanisms are not interchangeable: turning off the ceiling
assessment that triggers material escalation costs two of three
material-bottleneck tasks outright; turning off forced architecture exploration
costs one of six; turning off the three-model vote changes nothing---an
``unrealized premium'', since in this suite the governed agent never exercised
the molecular path, preferring cheaper bridges. Second, the governance is
model-robust: eight of eight under both a flagship and an economy LLM, the same
contracts attained through different designs. Third, the simulation world is
anchored to a physical artifact: the beginning-of-life $0.1$C discharge curve
of a production 21700 cell (LG M50T, \cite{lgm50dataset2024,
kirkaldy2024lgm50}) matches our calibrated parameter set
\cite{oregan2022electrochim} to within $0.45\%$ capacity and $20.6$ mV RMSE---a
validation of the simulator's calibration against a production cell, as
distinct from any validation of the designs themselves.

The paper makes four contributions. (i) We formulate agent-mediated design as
a governance problem: contract pre-registration, mechanized judgment, and
evidence-chain auditing are decision-process infrastructure rather than
simulation plumbing, and we measure the magnitudes they carry in controlled
data. (ii) We provide a three-condition benchmark under constant resources in
which the three-way comparison reveals that the prevailing ``LLM beats classic
optimization'' story \cite{chen2026battery} is incomplete: the two standard
baselines fail for different, specifiable reasons, and only protocol governance
fixes both. (iii) We contribute a mechanism decomposition of governance with
calibration: each component's marginal value is reported as measured,
including the one that turns out dormant in our suite. (iv) We measure decision
quality not as physics but as contract score---attainment rate, verifiable
evidence completeness, and resource (tokens and messages)---a transferable
yardstick for future agent-scale deployments in scientific and engineering
organizations.

The remainder of the paper is organized as follows. Section~2 reviews autonomy
in scientific and engineering design and situates governance in the delegation
literature. Section~3 specifies the protocol and its implementation. Section~4
defines the experimental design and the contract-verification machinery.
Section~5 reports results; Section~6 discusses and calibrates; Section~7
concludes with implications for the design and management of agent-mediated
decision systems.
